\chapter{Notebook Workflow and Repository Design}
\label{ch:notebooks}

The work is carried out in Jupyter notebooks. Notebooks are good for
exploration and poor at reproducibility. The design below keeps their
convenience and removes their usual failure modes.

\section{Principles}

\begin{enumerate}
  \item \emph{Notebooks drive; modules hold the code.} Every model, loss,
        dataset and metric lives in a tested Python package. Notebooks import
        it, set parameters, run and display. Nothing a result depends on is
        defined only inside a notebook cell.
  \item \emph{One notebook, two machines.} Each notebook detects whether it
        runs locally or on Kaggle and sets paths accordingly. The same file
        runs in both places without edits.
  \item \emph{Configuration as files.} Data, model, training and evaluation
        settings are YAML files under version control. A run is identified
        by the hash of its resolved configuration.
  \item \emph{Every run is logged.} Each run records five things:
        \begin{itemize}
          \item its seed;
          \item the git commit;
          \item the configuration hash;
          \item the library versions;
          \item its metrics, written as machine-readable JSON.
        \end{itemize}
  \item \emph{Results are generated, never typed.} The tables and figures of
        the paper are produced from the run logs by Notebook~15.
\end{enumerate}

\section{Repository layout}

\begin{Verbatim}
cffm-net/
  pyproject.toml            pinned dependencies
  configs/
    data/                   rgbt_tiny.yaml, llvip.yaml, m3fd.yaml, ...
    model/                  yolo26s_dual.yaml, cffm_net_s.yaml, ablations/
    train/                  stage_a.yaml, stage_b.yaml
    eval/                   robustness.yaml, latency.yaml
  src/cffm/
    data/                   paired and clip datasets, corruption simulator,
                            converters, alignment audit
    models/                 backbones, cmfm.py, scan.py (CUDA and PyTorch),
                            tsm.py, egomotion.py, monitor.py, registry.py
    train/                  stage_a.py, stage_b.py, losses.py, tbptt.py
    eval/                   coco_eval.py, size_bins.py, safit.py,
                            stream_metrics.py, track_eval.py, latency.py,
                            bootstrap.py
    deploy/                 export_trt.py, stream_runtime.py, capture.py,
                            webrtc_app.py
    utils/                  logging, seeding, checkpoint and resume
  notebooks/                00 to 15, described in Table 10.1
  tests/                    unit, property and regression tests
  scripts/                  kaggle_setup.sh, build_mamba_wheel.sh,
                            package_dataset.py
  runs/                     run logs (synchronised, not committed)
  paper/                    manuscript with generated tables
\end{Verbatim}

The model modules are registered with Ultralytics' model parser, so a model
is defined by a YAML file in the Ultralytics format and trained with its
trainer. A custom trainer subclass adds paired inputs, clip training and the
reliability loss.

\section{The notebook sequence}

Table~\ref{tab:notebooks} lists the sixteen notebooks, and
Figure~\ref{fig:notebooks} shows how they depend on one another.

\begin{table}[htbp]
\centering
\caption[The sixteen notebooks]{The sixteen notebooks. The location column
says where each normally runs: L is the local laptop, K is Kaggle.}
\label{tab:notebooks}
\small
\begin{tabularx}{\textwidth}{@{}P{0.65cm}P{3.6cm}LP{0.8cm}@{}}
\hd{No.} & \hd{Notebook} & \hd{Purpose and main output} & \hd{Runs} \\ \gap
00 & environment setup & install, GPU check, compile and verify \texttt{mamba\_ssm}, scan equivalence, smoke training & L, K \\
01 & data acquisition & downloads, request tracking, checksums, dataset cards & L \\
02 & data conversion & paired folders, YOLO labels, COCO JSON, ignore regions & L, K \\
03 & data audit & size histograms, density, class balance, alignment offsets & K \\
04 & single-modality baselines & YOLO26 on visible and thermal; measured throughput replaces estimates & K \\
05 & fusion baselines & two-stream fusion, CFT-style, Fusion-Mamba, COMO & K \\
06 & CMFM tests & gating, interleaving and reliability maps, visualised & L \\
07 & Stage A training & still-image \method{} with three seeds & K \\
08 & image ablations & A1 to A6 and A15 & K \\
09 & clips and TSM & clip loader, corruption simulator checks, constant-memory test & L, K \\
10 & Stage B training & \methodS{} and temporal baselines & K \\
11 & robustness evaluation & the protocol of Table~\ref{tab:robustness} & K \\
12 & tracking evaluation & ByteTrack, BoT-SORT, TrackTrack with TrackEval & K \\
13 & export and latency & TensorRT export, latency protocol, accuracy of the engine & L \\
14 & live demonstration & capture threads, stream runtime, FastRTC application & L \\
15 & results and tables & aggregation, bootstrap intervals, LaTeX tables and figures & L \\
\end{tabularx}
\end{table}

\begin{figure}[htbp]
\centering
\begin{tikzpicture}[
    nb/.style={base, text width=1.22cm, minimum height=1.0cm, font=\scriptsize},
    loc/.style={nb, rgb},
    kag/.style={nb, fuse},
    both/.style={nb, n},
  ]
  \def\dx{1.6}
  \node[both] (n00) at (0*\dx,0)    {\textbf{00}\\ setup};
  \node[loc]  (n01) at (1*\dx,0)    {\textbf{01}\\ acquire};
  \node[both] (n02) at (2*\dx,0)    {\textbf{02}\\ convert};
  \node[kag]  (n03) at (3*\dx,1.45) {\textbf{03}\\ audit};
  \node[kag]  (n04) at (3*\dx,0)    {\textbf{04}\\ single};
  \node[loc]  (n06) at (3*\dx,-1.45){\textbf{06}\\ CMFM tests};
  \node[both] (n09) at (3*\dx,-2.9) {\textbf{09}\\ clips, TSM};
  \node[kag]  (n05) at (4*\dx,0)    {\textbf{05}\\ fusion};
  \node[kag]  (n07) at (5*\dx,0)    {\textbf{07}\\ Stage A};
  \node[kag]  (n08) at (6*\dx,1.45) {\textbf{08}\\ ablations};
  \node[kag]  (n10) at (6*\dx,-1.45){\textbf{10}\\ Stage B};
  \node[kag]  (n11) at (7*\dx,0)    {\textbf{11}\\ robustness};
  \node[kag]  (n12) at (7*\dx,-1.45){\textbf{12}\\ tracking};
  \node[loc]  (n13) at (7*\dx,-2.9) {\textbf{13}\\ export};
  \node[loc]  (n15) at (8*\dx,1.45) {\textbf{15}\\ tables};
  \node[loc]  (n14) at (8*\dx,-2.9) {\textbf{14}\\ live demo};

  \draw[arr] (n00) -- (n01);
  \draw[arr] (n01) -- (n02);
  \draw[arr] (n02) -- (n04);
  \draw[arr] (n02.north) |- (n03.west);
  \draw[arr] (n02.south) |- (n09.west);
  \draw[arr] (n04) -- (n05);
  \draw[arr] (n05) -- (n07);
  \draw[arr] (n06.east) -| ($(n07.south)+(-0.25,0)$);
  \draw[darr] (n03.east) -| node[lbl, pos=0.25, above=0.5pt] {informs} ($(n07.north)+(-0.25,0)$);
  \draw[arr] (n07.north east) ++(-0.1,0) |- (n08.west);
  \draw[arr] (n07.south east) ++(-0.1,0) |- (n10.west);
  \draw[arr] (n09.east) -| ($(n10.south)+(0.0,0)$);
  \draw[arr] (n10.east) -- ++(0.2,0) |- (n11.west);
  \draw[arr] (n10) -- (n12);
  \draw[arr] (n10.south east) ++(-0.15,0) |- (n13.west);
  \draw[arr] (n13) -- (n14);
  \draw[arr] (n08) -- (n15);
  \draw[arr] (n11.north east) ++(-0.1,0) |- ($(n15.south)+(0,0)$);
  \draw[darr] (n12.east) -| ($(n15.south)+(0.35,0)$);

  % legend
  \node[loc, text width=1.2cm, minimum height=0.55cm] (l1) at (0.35*\dx,-2.9) {local};
  \node[kag, text width=1.2cm, minimum height=0.55cm] (l2) at (1.35*\dx,-2.9) {Kaggle};
  \node[both, text width=1.2cm, minimum height=0.55cm] (l3) at (0.85*\dx,-2.2) {both};
  \node[lbl, anchor=north] at (0.85*\dx,-3.25) {where each notebook runs};
\end{tikzpicture}

\vspace{2mm}
\caption[Dependencies between the notebooks]{Dependencies between the
notebooks. Colour shows where each notebook normally runs. Every training
and evaluation notebook writes run logs, and Notebook~15 reads all of them.
Only the main dependencies into it are drawn.}
\label{fig:notebooks}
\end{figure}


\section{Testing}

Three kinds of test protect the results.
\begin{itemize}
  \item \emph{Unit tests} check shapes and numerical agreement. The most
        important check is that the CUDA and pure-PyTorch scans agree.
  \item \emph{Property tests} check the behaviour the method relies on.
        \begin{itemize}
          \item With $\Delta=0$ the state does not change.
          \item With zero reliability for a modality, its tokens do not alter
                the state.
          \item Image mode equals a single step from an empty state.
          \item Memory use of the TSM does not grow with sequence length.
        \end{itemize}
  \item \emph{Regression tests} run a tiny model on a fixed fixture of
        frames. They fail if its outputs change unexpectedly.
\end{itemize}
The tests run in Notebook~00 on both machines and before every training
run.

\section{Generating results}
\label{sec:nb-results}

Notebook~15 is the only route from experiments to the manuscript.
\begin{enumerate}
  \item It reads every run log and groups runs by configuration.
  \item It computes means, standard deviations and sequence-level bootstrap
        intervals.
  \item It applies the multiple-comparison correction.
  \item It writes LaTeX tables and plot data into the paper's directory.
\end{enumerate}
A number in the paper can therefore always be traced to the runs that
produced it. A changed run changes the paper on the next build, and no table
is edited by hand.
